% SPDX-License-Identifier: CC-BY-SA-4.0
% Author: Matthieu Perrin
% Part: <Nom de la partie>
% Section: <Nom de la section>
% Sub-section: <Nom de la sous-section>  % (facultatif, laisser vide si non utilisé)
% Frame: <Titre de la slide>

\begingroup

\newcommand\Tpolyo[2][black]{
  \draw[#1, fill=#1!30] #2 ++(0.1,0.1) rectangle +(.9,.9);
  \draw[#1, fill=#1!30] #2 ++(1.1,0.1) rectangle +(.9,.9);
  \draw[#1, fill=#1!30] #2 ++(2.1,0.1) rectangle +(.9,.9);
  \draw[#1, fill=#1!30] #2 ++(1.1,1.1) rectangle +(.9,.9);
}

\begin{frame}{Problème de pavage du plan}

  \on[text, top=-3mm]{
    \Probleme{polyomino tiling}{
      Un ensemble fini de \structure{polyominos}
      \begin{itemize}
      \item ils sont composés de carrés unitaires collés
      \item ils peuvent être répétés et tournés
      \end{itemize}
    }{
      Peut-on recouvrir entièrement le \structure{plan infini} ?
      \begin{itemize}
        \item chaque point du plan doit être couvert
        \item les polyominos ne peuvent pas se chevaucher
      \end{itemize}
    }
  }

  \onBlock[y=-5mm]{Théorème (admis)}{
    \textsc{polyomino tiling} est indécidable.
  }
    
  \onExampleBlock[y=-26mm]{Exemples}{
    \begin{itemize}
    \item\vspace{-1mm}
      $\left\{
      \begin{tikzpicture}[size=3mm, baseline=(base)]
        \coordinate (base) at (0,0.75);
        \Tpolyo[structure]{(0,0)}
      \end{tikzpicture}
      \right\}$
      est une instance positive
    \item
      $\left\{
      \begin{tikzpicture}[size=3mm, baseline=(base)]
        \coordinate (base) at (0,1.25);
        \draw[example, fill=example!30] (0,0) ++(0.1,0.1) rectangle +(.9,.9);
        \draw[example, fill=example!30] (1,0) ++(0.1,0.1) rectangle +(.9,.9);
        \draw[example, fill=example!30] (2,0) ++(0.1,0.1) rectangle +(.9,.9);
        \draw[example, fill=example!30] (3,0) ++(0.1,0.1) rectangle +(.9,.9);
        \draw[example, fill=example!30] (4,0) ++(0.1,0.1) rectangle +(.9,.9);
        \draw[example, fill=example!30] (0,1) ++(0.1,0.1) rectangle +(.9,.9);
        \draw[example, fill=example!30] (4,1) ++(0.1,0.1) rectangle +(.9,.9);
        \draw[example, fill=example!30] (0,2) ++(0.1,0.1) rectangle +(.9,.9);
        \draw[example, fill=example!30] (1,2) ++(0.1,0.1) rectangle +(.9,.9);
        \draw[example, fill=example!30] (3,2) ++(0.1,0.1) rectangle +(.9,.9);
        \draw[example, fill=example!30] (4,2) ++(0.1,0.1) rectangle +(.9,.9);
      \end{tikzpicture}
      \right\}$
      est une instance négative
    \end{itemize}
  }

  \on[bottom=1mm, x=30mm]{
    \begin{tikzpicture}[size=3mm]
      \begin{scope}
        \clip (1.05,1.05) rectangle (12.05,12.05);
        \Tpolyo[example]   {(2 ,11)}
        \Tpolyo[example]   {(6 ,11)}
        \Tpolyo[example]   {(10,11)}
        \Tpolyo[structure] {(0 ,10)}
        \Tpolyo[structure] {(4 ,10)}
        \Tpolyo[structure] {(8 ,10)}
        \Tpolyo[alert]     {(2 ,9)}
        \Tpolyo[alert]     {(6 ,9)}
        \Tpolyo[alert]     {(10,9)}
        \Tpolyo[example]   {(0 ,8)}
        \Tpolyo[example]   {(4 ,8)}
        \Tpolyo[example]   {(8 ,8)}
        \Tpolyo[structure] {(2 ,7)}
        \Tpolyo[structure] {(6 ,7)}
        \Tpolyo[structure] {(10,7)}
        \Tpolyo[alert]     {(0 ,6)}
        \Tpolyo[alert]     {(4 ,6)}
        \Tpolyo[alert]     {(8 ,6)}
        \Tpolyo[example]   {(2 ,5)}
        \Tpolyo[example]   {(6 ,5)}
        \Tpolyo[example]   {(10,5)}
        \Tpolyo[structure] {(0 ,4)}
        \Tpolyo[structure] {(4 ,4)}
        \Tpolyo[structure] {(8 ,4)}
        \Tpolyo[alert]     {(2 ,3)}
        \Tpolyo[alert]     {(6 ,3)}
        \Tpolyo[alert]     {(10,3)}
        \Tpolyo[example]   {(0 ,2)}
        \Tpolyo[example]   {(4 ,2)}
        \Tpolyo[example]   {(8 ,2)}
        \Tpolyo[structure] {(2 ,1)}
        \Tpolyo[structure] {(6 ,1)}
        \Tpolyo[structure] {(10,1)}
        \Tpolyo[alert]     {(0 ,0)}
        \Tpolyo[alert]     {(4 ,0)}
        \Tpolyo[alert]     {(8 ,0)}
      \end{scope}
    \end{tikzpicture}
  }

  \footnoteref{N. Ollinger. \emph{Tiling the plane with a fixed number of polyominoes.} LATA (2009)}
  
\end{frame}

\endgroup
